% GENERATED FILE. Edit canonical lessons, then run npm run generate:books.
% canonical-source-hash: fnv1a64:6d07e0d8a49ad68b
% canonical-lessons: ES-C311-tienda, ES-C311-cuenta, ES-C311-precio, ES-C311-bolsa, ES-C311-asiento

\chapter{La Tienda Y La Cuenta --- Shop, Bill, Price, Bag, Seat}
\label{ch:la-tienda-y-la-cuenta}
\hlchaptermodality{311}

\begin{chapteropening}
\textbf{By the end of this chapter:} \emph{I can carry out a transaction in Spanish: find the shop, ask the price, ask for the bill, ask for a bag and ask for a seat.}

Complete ES-C311-asiento: walk in, ask the price, ask for the bill, take the bag and ask for a seat.
\end{chapteropening}

\section[la tienda]{la tienda — a shop that is really a stretched cloth}

\label{lesson:ES-C311-tienda}



\begin{quote}
Two from last time. What is \emph{la luz}? (The light.) And is English \emph{bath} a relative of \emph{baño}? (No — the real one is \emph{balneology}.)
\end{quote}

\begin{sounds}
\begin{itemize}
\raggedright
  \item \emph{TIEN-da}, stress on the first part. The \emph{ie} is one glide, not two beats. $\to$ pronunciation reference
\end{itemize}
\end{sounds}

\begin{cousinweb}
\textbf{la tienda} — \textquotedblleft{}the shop.\textquotedblright{}

Latin \textbf{tendere} meant \textbf{to stretch}, and a \emph{tienda} is named for the cloth that was stretched over it. Before a shop was a building it was an awning in a market: a piece of canvas pulled tight over poles, with goods laid out underneath. Spanish still uses \emph{tienda} for the camping kind as well, which is the older sense showing through.

English holds the same root in a wide spread. A \textbf{tent} is the plain one — the stretched cloth itself. Something \textbf{tense} is pulled tight. To \textbf{extend} is to stretch out; to \textbf{intend} is to stretch your mind \emph{toward} a thing; to \textbf{attend} is to stretch toward it and stay. And a \textbf{tendon} is the cord in your ankle that is stretched taut every time you take a step.

So a shop, a tent, an intention and an ankle are all one act: something pulled tight.
\end{cousinweb}

\begin{grammarlens}[title={What you've built}]
The first word of the run where money changes hands.
\end{grammarlens}

\subsection*{Your turn}
Say these aloud:

\begin{itemize}
\raggedright
  \item \emph{la tienda}
  \item \emph{¿dónde está la tienda?}
  \item \emph{disculpe, ¿dónde está la tienda?}
  \item \emph{la luz}, then \emph{el baño}, then \emph{la tienda}
\end{itemize}

\begin{tcolorbox}[breakable,colback=teal!4,colframe=teal!35!black,title={Before you move on}]
What did \emph{tendere} mean? (To stretch.) Which English word is the plain stretched cloth? (\emph{Tent}.) And what is \emph{el baño}? (The bathroom.)
\end{tcolorbox}
\section[la cuenta]{la cuenta — one Latin verb, arriving in English twice}

\label{lesson:ES-C311-cuenta}



\begin{quote}
Two from last time. What is \emph{el baño}? (The bathroom.) And what was a \emph{tienda} before it was a shop? (A stretched cloth.)
\end{quote}

\begin{sounds}
\begin{itemize}
\raggedright
  \item \emph{CUEN-ta}, stress on the first part. The \emph{c} is hard, and \emph{ue} is one glide. $\to$ pronunciation reference
\end{itemize}
\end{sounds}

\begin{cousinweb}
\textbf{la cuenta} — \textquotedblleft{}the bill.\textquotedblright{}

Latin \textbf{computare} comes apart into \emph{com-}, \textquotedblleft{}together,\textquotedblright{} and \emph{putare}, \textquotedblleft{}to reckon.\textquotedblright{} To \emph{computare} was to \textbf{reckon things up together} — to bring separate amounts into one number.

Here is the good part. English received this word \textbf{twice}, by two different roads, and did not notice they were the same.

Straight from Latin, late and in one piece, it arrived as \textbf{compute} — and from there, \textbf{computer}. Down through spoken Latin into French and then into English, worn smaller at every step, the same verb arrived as \textbf{count}, and with the prefix still attached, as \textbf{account} and \textbf{accountant}.

\emph{Compute} and \emph{count} are one word wearing different amounts of wear. A \textbf{computer} and an \textbf{accountant} have the same job description in Latin.

Spanish kept the road that wears words down, which is how \emph{computare} ends up as \emph{cuenta}. Ask for it after a meal: \emph{la cuenta, por favor}.
\end{cousinweb}

\begin{grammarlens}[title={What you've built}]
A shop and a bill: the two halves of any transaction.
\end{grammarlens}

\subsection*{Your turn}
Say these aloud:

\begin{itemize}
\raggedright
  \item \emph{la cuenta}
  \item \emph{la cuenta, por favor}
  \item \emph{disculpe — la cuenta, por favor}
  \item \emph{el baño}, then \emph{la tienda}, then \emph{la cuenta}
\end{itemize}

\begin{tcolorbox}[breakable,colback=teal!4,colframe=teal!35!black,title={Before you move on}]
Take \emph{computare} apart. (\emph{Com-} together, \emph{putare} to reckon.) Which two English words are it, one worn and one not? (\emph{Count} and \emph{compute}.) And what is \emph{la tienda}? (The shop.)
\end{tcolorbox}
\section[el precio]{el precio — why praising somebody is pricing them}

\label{lesson:ES-C311-precio}



\begin{quote}
Two from last time. What was a \emph{tienda} made of? (Stretched cloth.) And which two English words are both \emph{computare}? (\emph{Count} and \emph{compute}.)
\end{quote}

\begin{sounds}
\begin{itemize}
\raggedright
  \item \emph{PRE-cio}, stress on the first part. The \emph{pr} runs together with no vowel between, and the \emph{-cio} ending is one glide. $\to$ pronunciation reference
\end{itemize}
\end{sounds}

\begin{cousinweb}
\textbf{el precio} — \textquotedblleft{}the price.\textquotedblright{}

Latin \textbf{pretium} was \textbf{what a thing is worth}. Not the number on a tag: the worth itself, which the number is only a report of.

English carries the family in every register at once. \textbf{Price} is the plain transaction. A \textbf{prize} is a thing of worth handed over. Something \textbf{precious} has worth in it. To \textbf{appraise} is to set a worth on a thing, and to \textbf{appreciate} is to \emph{raise} its worth — which is why a house appreciates and why appreciating a friend is a real compliment rather than a soft one.

And then the one that surprises people: \textbf{praise}. To praise somebody is \emph{pretiare}, to put a value on them out loud. Praise is not decoration. It is an appraisal, said where the person can hear it.
\end{cousinweb}

\begin{grammarlens}[title={What you've built}]
Shop, bill, price. You can now ask what a transaction costs before it happens rather than after.
\end{grammarlens}

\subsection*{Your turn}
Say these aloud:

\begin{itemize}
\raggedright
  \item \emph{el precio}
  \item \emph{el precio, por favor}
  \item \emph{disculpe — el precio, por favor}
  \item \emph{la tienda}, then \emph{la cuenta}, then \emph{el precio}
\end{itemize}

\begin{tcolorbox}[breakable,colback=teal!4,colframe=teal!35!black,title={Before you move on}]
What did \emph{pretium} mean? (What a thing is worth.) What are you doing to somebody when you \emph{praise} them, going by the root? (Putting a value on them.) And what is \emph{la cuenta}? (The bill.)
\end{tcolorbox}
\section[la bolsa]{la bolsa — a bag made of skin, and a word made of it too}

\label{lesson:ES-C311-bolsa}



\begin{quote}
Two from last time. What is \emph{la cuenta}? (The bill.) And what are you really doing when you \emph{praise} somebody? (Putting a value on them — \emph{pretium}.)
\end{quote}

\begin{sounds}
\begin{itemize}
\raggedright
  \item \emph{BOL-sa}, stress on the first part. The \emph{b} between vowels is soft, with the lips barely closing. $\to$ pronunciation reference
\end{itemize}
\end{sounds}

\begin{cousinweb}
\textbf{la bolsa} — \textquotedblleft{}the bag.\textquotedblright{}

Latin \textbf{bursa} was a leather bag, and Latin had taken it from Greek \emph{byrsa}, a stripped \textbf{hide}. Before a bag was cloth or paper it was skin, because skin was what came off the animal already waterproof and already the right size.

English took \emph{bursa} through French and made \textbf{purse} of it. Then it took the Latin form again, unworn, for the offices that handle money: a \textbf{bursar} keeps a college's purse, and a \textbf{bursary} is what the purse pays out. To \textbf{disburse} is to pay \emph{out of} the bag; to \textbf{reimburse} is to put back \emph{into} it. Every one of those words is a hide.

Spanish uses \emph{bolsa} for the shopping bag you are handed at a counter, and — in exactly the same logic English used for \emph{purse} — for the stock exchange, the place where the money bag of a whole country is opened.
\end{cousinweb}

\begin{grammarlens}[title={What you've built}]
You can be given the bag as well as the bill now.
\end{grammarlens}

\subsection*{Your turn}
Say these aloud:

\begin{itemize}
\raggedright
  \item \emph{la bolsa}
  \item \emph{la bolsa, por favor}
  \item \emph{disculpe — la bolsa, por favor}
  \item \emph{la cuenta}, then \emph{el precio}, then \emph{la bolsa}
\end{itemize}

\begin{tcolorbox}[breakable,colback=teal!4,colframe=teal!35!black,title={Before you move on}]
What was a \emph{byrsa}? (A stripped hide.) What does a \emph{bursar} look after? (The purse.) And what is \emph{el precio}? (The price.)
\end{tcolorbox}
\section[el asiento]{el asiento — everything that settles}

\label{lesson:ES-C311-asiento}



\begin{quote}
Two from last time. What is \emph{el precio}? (The price.) And what was a \emph{bolsa} made of, at the start? (A hide.)
\end{quote}

\begin{sounds}
\begin{itemize}
\raggedright
  \item \emph{a-SIEN-to}, stress on the middle. Four letters do the middle syllable and it is still one beat. $\to$ pronunciation reference
\end{itemize}
\end{sounds}

\begin{cousinweb}
\textbf{el asiento} — \textquotedblleft{}the seat.\textquotedblright{}

Under it is Latin \textbf{sedere}, \textbf{to sit}, with \emph{ad-} in front: to sit down \emph{to} something. Spanish built the verb \emph{asentar} and then this noun off it.

\emph{Sedere} has settled into English everywhere. Work that keeps you in a chair is \textbf{sedentary}. What settles to the bottom of a river is \textbf{sediment}. A \textbf{session} is a sitting — of a court, a parliament, a band. To \textbf{reside} is to sit \emph{back}, to stay put in a place. A bishop's \textbf{see} is his seat, which is why the Vatican is the Holy See rather than the Holy Anything Else.

And \textbf{assess}: \emph{ad-} plus \emph{sedere} again, the same build as \emph{asiento}. An assessor was the official who \textbf{sat beside} the judge and set the tax. To assess is not to look from above. It is to sit next to a thing and reckon it.

\emph{Un asiento, por favor} is how you ask for a place on a bus.
\end{cousinweb}

\begin{grammarlens}[title={What you've built}]
Five words for a transaction: the shop, the bill, the price, the bag, the seat. Between them they cover walking in, paying, carrying it out and sitting down.

Next comes what you are carrying while you do it.
\end{grammarlens}

\subsection*{Your turn}
Say these aloud:

\begin{itemize}
\raggedright
  \item \emph{el asiento}
  \item \emph{un asiento, por favor}
  \item \emph{disculpe — un asiento, por favor}
  \item \emph{el precio}, then \emph{la bolsa}, then \emph{el asiento}
  \item \emph{gracias}
\end{itemize}

\begin{tcolorbox}[breakable,colback=teal!4,colframe=teal!35!black,title={Before you move on}]
What does an assessor do, literally? (Sits beside.) Which English word means work done in a chair? (\emph{Sedentary}.) And what is \emph{la bolsa}? (The bag.)
\end{tcolorbox}
